\section{Discussion}

The path from a distillation loss to task performance passes through optimizer history, numerical precision, and an evolving training trajectory. In our comparisons, Adam aligns directionally different gradients, and BF16 rounding hides widespread FP32 changes. Broader vocabulary supervision improves gradient approximation and produces task-dependent learning gains, including a mathematics--science trade-off and changing advantages over time. These findings give complementary roles to the measurements: gradients describe the immediate training signal, parameter changes show its interaction with the optimizer and precision, and per-task learning curves reveal the resulting capability gains and regressions.
